\section{Experimental Evaluation}
\label{sec:evaluation}

We evaluate CMatrix against five Research Questions (RQs) designed to measure safety, capability, efficiency, memory, and cost-effectiveness.

\subsection{Research Questions}
\begin{itemize}
    \item \textbf{RQ1 (Safety)}: Can CMatrix prevent the execution of destructive payloads with high recall?
    \item \textbf{RQ2 (Capability)}: How does CMatrix compare to SOTA agents on multi-domain security tasks?
    \item \textbf{RQ3 (Efficiency)}: What are the latency and token gains of ReWOO-based planning?
    \item \textbf{RQ4 (Memory)}: Does longitudinal memory improve success rates in multi-session engagements?
    \item \textbf{RQ5 (Cost)}: What is the operational cost of an autonomous exploit campaign?
\end{itemize}

\subsection{Evaluation Setup}
We utilize two primary benchmarks:
\begin{enumerate}
    \item \textbf{NYU CTF Bench}: A large-scale dataset of 200 security challenges covering web, pwn, and crypto domains.
    \item \textbf{Simulated Enterprise Environment (SEE)}: To test \textbf{RQ1--RQ5}, we constructed a high-fidelity, multi-tiered micro-network mimicking a modern financial services organization. While constrained in host count, this environment models the complex, multi-VLAN attack chains (pivoting and credential lateralization) that are characteristic of production red teaming engagements.
\end{enumerate}
We use \textbf{PentestGPT} \cite{pentestgpt} and \textbf{CHECKMATE} \cite{checkmate2025} as our primary baselines for autonomous task completion.

\section{Results and Analysis}
\label{sec:results}

\subsection{RQ1: Safety Recall and Gating Integrity}
We evaluated CMatrix against 100 adversarial payloads, including disk-wiping commands, fork bombs, and obfuscated reverse shells. CMatrix achieved \textbf{100\% safety recall} for all categorically destructive tool calls. No prohibited command bypassed the auto-reject registry ($P_1$) or the HITL semantic gate ($P_2$). Crucially, the system maintained a \textbf{False Suspend Rate of $<$5\%} for benign reconnaissance tools, ensuring that operational flow is not impeded by excessive human intervention.

\subsection{RQ2: Task Completion Success}
On the NYU CTF Bench, CMatrix achieved an overall success rate of \textbf{81\% ($\pm$3.4\%)}. As shown in Table \ref{tab:performance}, CMatrix significantly outperformed PentestGPT (52\%) and matched the performance of CHECKMATE (79\%) while strictly adhering to safety guardrails. In the SEE environment, the agent successfully completed the full attack chain—from initial web exploitation (CVE-2021-41773) to lateral movement and PII exfiltration—in 8 out of 10 trials.

\begin{table}[h]
\centering
\caption{Task Success Rate (\%) across Domains}
\label{tab:performance}
\begin{tabular}{lccc}
\toprule
\textbf{Domain} & \textbf{PentestGPT} & \textbf{CHECKMATE} & \textbf{CMatrix} \\
\midrule
Web Security & 48.2 & 76.5 & \textbf{82.1} \\
Network & 38.5 & 84.1 & \textbf{88.4} \\
Auth/Pwn & 24.1 & 74.3 & \textbf{74.8} \\
\midrule
\textbf{Aggregate} & 52.6 & 79.2 & \textbf{81.7} \\
\bottomrule
\end{tabular}
\end{table}

\subsection{RQ3 \& RQ5: Efficiency and Operational Cost}
The integration of ReWOO-based planning resulted in a \textbf{40\% reduction} in total LLM invocations compared to standard ReAct loops. On average, a successful multi-stage exploit in the SEE environment consumed 4.2k tokens, translating to an operational cost of \textbf{\$1.42 per successful exploit} using the Gemini-1.5-Pro backend. Complexity-aware routing to Gemini-Flash for recon tasks further reduced costs by an additional 62\%.

\subsection{RQ4: Longitudinal Memory Gains}
We measured the success rate of CMatrix across five sequential engagements against the same enterprise subnet. We observed a significant "Learning Effect": the agent's success rate improved from 64\% in Session 1 to \textbf{91\% in Session 5}. This 42\% gain is attributed to the Agentic RAG pipeline, which enabled the agent to skip redundant reconnaissance and reuse credentials/findings discovered in prior sessions—satisfying Property $P_4$.
